L'operateur de collisions peut etre exprime sous la forme de la divergenced'un
flux comme suit: \[
C(f_{0})=-\nabla_{p}\cdot S^{\text{C}}=-\frac{1}{p^{2}}\frac{\partial}{\partial p}\left(p^{2}S_{p}^{\text{C}}\right)+\frac{1}{p}\frac{\partial}{\partial\xi}\left(\sqrt{1-\xi^{2}}S_{\xi}^{\text{C}}\right)\]
avec \[
\left\{ \begin{array}{l}
S_{p}^{\text{C}}=-\left[A(p)\frac{\partial f_{0}}{\partial p}+F(p)f_{0}\right]\\
S_{\xi}^{\text{C}}=\frac{B_{t}(p)}{p}\sqrt{1-\xi^{2}}\frac{\partial f_{0}}{\partial\xi}\end{array}\right.\]
ou \[
S^{\text{C}}=-\overline{\overline{D}}_{p}^{\text{C}}\cdot\nabla_{p}f_{0}+\overline{F}_{p}^{\text{C}}f_{0}\]
avec \[
\overline{\overline{D}}_{p}^{\text{C}}\rightarrow\left\{ \begin{array}{l}
D_{pp}^{\text{C}}=A(p)\\
D_{p\xi}^{\text{C}}=0\\
D_{\xi p}^{\text{C}}=0\\
D_{\xi\xi}^{\text{C}}=B_{t}(p)\end{array}\right.\]
et \[
\overline{F}_{p}^{\text{C}}\rightarrow\left\{ \begin{array}{l}
F_{p}^{\text{C}}=-F(p)\\
F_{\xi}^{\text{C}}=0\end{array}\right.\]


A partir de ces expressions il est facile de determiner la moyenne
des coefficients sur la periode de rebond, sur la base du formalisme
developpe au paragrapheprecedent.On obtient ainsi de fa\c{c}on immediate
\[
\left\{ \begin{array}{l}
D_{pp}^{\text{C}(0)}=A(p)\\
D_{p\xi}^{\text{C}(0)}=0\\
D_{\xi p}^{\text{C}(0)}=0\\
D_{\xi\xi}^{\text{C}(0)}=B_{t}(p)\left\{ \frac{\xi^{2}}{\Psi\xi_{0}^{2}}\right\} \end{array}\right.\]
et \[
\left\{ \begin{array}{l}
F_{p}^{\text{C}(0)}=-F(p)\\
F_{\xi}^{\text{C}(0)}=0\end{array}\right.\]


Le coefficient $\left\{ \xi^{2}/\Psi\xi_{0}^{2}\right\} $peut etre
reecrit la forme usuelle \begin{eqnarray*}
\left\{ \frac{\xi^{2}}{\Psi\xi_{0}^{2}}\right\}  & = & \left\{ \frac{1-\Psi\left(1-\xi_{0}^{2}\right)}{\Psi\xi_{0}^{2}}\right\} \\
 & = & \left\{ 1-\frac{1}{\xi_{0}^{2}}\left(1-\frac{1}{\Psi}\right)\right\} \\
 & = & 1-\frac{\Delta_{b}}{\xi_{0}^{2}}\end{eqnarray*}
avec \[
\Delta_{b}\equiv\left\{ 1-\frac{1}{\Psi}\right\} \]


Finalement, on obtient \[
\left\{ \begin{array}{l}
D_{pp}^{\text{C}(0)}=A(p)\\
D_{p\xi}^{\text{C}(0)}=0\\
D_{\xi p}^{\text{C}(0)}=0\\
D_{\xi\xi}^{\text{C}(0)}=\left(1-\frac{\Delta_{b}}{\xi_{0}^{2}}\right)B_{t}(p)\end{array}\right.\]
et \[
\left\{ \begin{array}{l}
F_{p}^{\text{C}(0)}=-F(p)\\
F_{\xi}^{\text{C}(0)}=0\end{array}\right.\]


Sous forme discrete, les coefficients qui interviennent dans la matrice$\overline{\overline{M}}^{\left(0\right)}$sont
\[
\overline{\overline{D}}_{p}^{\text{C}(0)}\rightarrow\left\{ \begin{array}{l}
D_{pp,i+1,j+1/2}^{\text{C}\left(0\right)}=A_{i+1}\\
D_{pp,i,j+1/2}^{\text{C}\left(0\right)}=A_{i}\\
D_{p\xi,i+1,j+1/2}^{\text{C}\left(0\right)}=0\\
D_{\xi p,i+1/2,j+1}^{\text{C}\left(0\right)}=0\\
D_{p\xi,i+1/2,j+1/2}^{\text{C}\left(0\right)}=0\\
D_{\xi p,i+1/2,j+1/2}^{\text{C}\left(0\right)}=0\\
D_{p\xi,i,j+1/2}^{\text{C}\left(0\right)}=0\\
D_{\xi p,i+1/2,j}^{\text{C}\left(0\right)}=0\\
D_{\xi\xi,i+1/2,j+1}^{\text{C}\left(0\right)}=\left(1-\Delta_{bj+1}/\xi_{0j+1}^{2}\right)B_{t,i+1/2}\\
D_{\xi\xi,i+1/2,j}^{\text{C}\left(0\right)}=\left(1-\Delta_{bj}/\xi_{0j}^{2}\right)B_{t,i+1/2}\end{array}\right.\]
et \[
\overline{F}_{p}^{\text{C}(0)}\rightarrow\left\{ \begin{array}{l}
F_{p,i+1,j+1/2}^{\text{C}(0)}=-F_{i+1}\\
F_{p,i,j+1/2}^{\text{C}(0)}=-F_{i}\\
F_{\xi,i+1/2,j+1}^{\text{C}(0)}=0\\
F_{\xi,i+1/2,j}^{\text{C}(0)}=0\end{array}\right.\]

